\ifdefined\gkver\else\def\gkver{student}\fi
\documentclass[\gkver]{gaokaozhenti}
\begin{document}

\chapter{2013年江苏}

\begin{enumerate}
\item
%% number: 1
%% paperName: 2013年江苏高考第 1 题 3 分
%% typeId: 1
%% type: 单选题
%% chapter: 曲线运动
%% point: 开普勒定律
%% method: 无
%% score: 3
%% degree: 850
%% duplicateId: 0
%% body:
火星和木星沿各自的椭圆轨道绕太阳运行，根据开普勒行星运动定律可知\xzanswer{C}
\fourchoices[answer=C]
{太阳位于木星运行轨道的中心}
{火星和木星绕太阳运行速度的大小始终相等}
{火星与木星公转周期之比的平方等于它们轨道半长轴之比的立方}
{相同时间内，火星与太阳连线扫过的面积等于木星与太阳连线扫过的面积}
%% answer: C
%% memo:
\memoanswer{
本题原卷及材料中未提供详解，待补充。
}

\item
%% number: 2
%% paperName: 2013年江苏高考第 2 题 3 分
%% typeId: 1
%% type: 单选题
%% chapter: 曲线运动
%% point: 圆周运动的应用
%% method: 无
%% score: 3
%% degree: 650
%% duplicateId: 0
%% body:
如图所示，“旋转秋千”中的两个座椅A、B质量相等，通过相同长度的缆绳悬挂在旋转圆盘上。不考虑空气阻力的影响，当旋转圆盘绕竖直的中心轴匀速转动时，下列说法正确的是\xzanswer{D}
\onepicture[w=3.5cm]{figs/02.png}
\fourchoices[answer=D]
{A的速度比B的大}
{A与B的向心加速度大小相等}
{悬挂A、B的缆绳与竖直方向的夹角相等}
{悬挂A的缆绳所受的拉力比悬挂B的小}
%% answer: D
%% memo:
\memoanswer{
本题原卷及材料中未提供详解，待补充。
}

\item
%% number: 3
%% paperName: 2013年江苏高考第 3 题 3 分
%% typeId: 1
%% type: 单选题
%% chapter: 电场
%% point: 电场叠加
%% method: 无
%% score: 3
%% degree: 650
%% duplicateId: 0
%% body:
下列选项中的各$\frac{1}{4}$圆环大小相同，所带电荷量已在图中标出，且电荷均匀分布，各$\frac{1}{4}$圆环间彼此绝缘。坐标原点$O$处电场强度最大的是\xzanswer{B}
\fourchoices[answer=B, ispicture=true, h=2.0cm]
{figs/03a.png}
{figs/03b.png}
{figs/03c.png}
{figs/03d.png}
%% answer: B
%% memo:
\memoanswer{
本题原卷及材料中未提供详解，待补充。
}

\item
%% number: 4
%% paperName: 2013年江苏高考第 4 题 3 分
%% typeId: 1
%% type: 单选题
%% chapter: 电路
%% point: 动态分析
%% method: 无
%% score: 3
%% degree: 650
%% duplicateId: 0
%% body:
在输液时，药液有时会从针口流出体外，为了及时发现，设计了一种报警装置，电路如图所示。M是贴在针口处的传感器，接触到药液时其电阻$R_{M}$发生变化，导致S两端电压$U$增大，装置发出警报，此时\xzanswer{C}
\onepicture[w=3.2cm]{figs/04.png}
\fourchoices[answer=C]
{$R_{M}$变大，且$R$越大，$U$增大越明显}
{$R_{M}$变大，且$R$越小，$U$增大越明显}
{$R_{M}$变小，且$R$越大，$U$增大越明显}
{$R_{M}$变小，且$R$越小，$U$增大越明显}
%% answer: C
%% memo:
\memoanswer{
接触到药液时其电阻$R_{M}$发生变化时，警报器S两端电压$U$增大，可知干路电流增大$\to$外电阻减小$\to R_{M}$减小；

可用\textbf{极限法}分析$U$增大灵敏度和什么因素有关：假设$R$很小，甚至为零，则传感器部分的电路被短路，故传感器$R_{M}$的大小变化对S的电压就无影响，而$R$越大，甚至接近于断路，那么传感器的变化就直接影响了S的电压的变化。所以，$R$越大，$U$增大越明显。

正确选项为C。
}

\item
%% number: 5
%% paperName: 2013年江苏高考第 5 题 3 分
%% typeId: 1
%% type: 单选题
%% chapter: 运动和力
%% point: 动量守恒定律
%% method: 无
%% score: 3
%% degree: 650
%% duplicateId: 0
%% body:
水平面上，一白球与一静止的灰球碰撞，两球质量相等。碰撞过程的频闪照片如图所示，据此可推断，碰撞过程中系统损失的动能约占碰撞前动能的\xzanswer{A}
\onepicture[w=4.5cm]{figs/05.png}
\fourchoices[answer=A]
{$30\%$}
{$50\%$}
{$70\%$}
{$90\%$}
%% answer: A
%% memo:
\memoanswer{
本题原卷及材料中未提供详解，待补充。
}

\item
%% number: 6
%% paperName: 2013年江苏高考第 6 题 4 分
%% typeId: 2
%% type: 多选题
%% chapter: 电场
%% point: 电场线
%% method: 无
%% score: 4
%% degree: 650
%% duplicateId: 0
%% body:
将一电荷量为$+Q$的小球放在不带电的金属球附近，所形成的电场线分布如图所示，金属球表面的电势处处相等。$a$、$b$为电场中的两点，则\xzanswer{ABD}
\onepicture[w=5.0cm]{figs/06.png}
\fourchoices[answer=ABD]
{$a$点的电场强度比$b$点的大}
{$a$点的电势比$b$点的高}
{检验电荷$-q$在$a$点的电势能比在$b$点的大}
{将检验电荷$-q$从$a$点移到$b$点的过程中，电场力做负功}
%% answer: ABD
%% memo:
\memoanswer{
本题原卷及材料中未提供详解，待补充。
}

\item
%% number: 7
%% paperName: 2013年江苏高考第 7 题 4 分
%% typeId: 2
%% type: 多选题
%% chapter: 曲线运动
%% point: 斜抛运动
%% method: 无
%% score: 4
%% degree: 650
%% duplicateId: 0
%% body:
如图所示，从地面上同一位置抛出两小球A、B，分别落在地面上的M、N点，两球运动的最大高度相同。空气阻力不计，则\xzanswer{CD}
\onepicture[w=5.5cm]{figs/07.png}
\fourchoices[answer=CD]
{B的加速度比A的大}
{B的飞行时间比A的长}
{B在最高点的速度比A在最高点的大}
{B在落地时的速度比A在落地时的大}
%% answer: CD
%% memo:
\memoanswer{
本题原卷及材料中未提供详解，待补充。
}

\item
%% number: 8
%% paperName: 2013年江苏高考第 8 题 4 分
%% typeId: 2
%% type: 多选题
%% chapter: 交流电
%% point: LC振荡回路
%% method: 无
%% score: 4
%% degree: 650
%% duplicateId: 0
%% body:
如图所示，理想变压器原线圈接有交流电源，当副线圈上的滑片P处于图示位置时，灯泡L能发光。要使灯泡变亮，可以采取的方法有\xzanswer{BC}
\onepicture[w=4.5cm]{figs/08.png}
\fourchoices[answer=BC]
{向下滑动P}
{增大交流电源的电压}
{增大交流电源的频率}
{减小电容器C的电容}
%% answer: BC
%% memo:
\memoanswer{
本题原卷及材料中未提供详解，待补充。
}

\item
%% number: 9
%% paperName: 2013年江苏高考第 9 题 4 分
%% typeId: 2
%% type: 多选题
%% chapter: 功和能
%% point: 动能定理
%% method: 无
%% score: 4
%% degree: 450
%% duplicateId: 0
%% body:
如图所示，水平桌面上的轻质弹簧一端固定，另一端与小物块相连。弹簧处于自然长度时物块位于O点（图中未标出）。物块的质量为$m$，AB$=a$，物块与桌面间的动摩擦因数为$\mu$。现用水平向右的力将物块从O点拉至A点，拉力做的功为$W$。撤去拉力后物块由静止向左运动，经O点到达B点时速度为零。重力加速度为$g$。则上述过程中\xzanswer{BC}
\onepicture[w=5.5cm]{figs/09.png}
\fourchoices[answer=BC]
{物块在A点时，弹簧的弹性势能等于$W-\frac{1}{2}\mu mga$}
{物块在B点时，弹簧的弹性势能小于$W-\frac{3}{2}\mu mga$}
{经O点时，物块的动能小于$W-\mu mga$}
{物块动能最大时弹簧的弹性势能小于物块在B点时弹簧的弹性势能}
%% answer: BC
%% memo:
\memoanswer{
本题原卷及材料中未提供详解，待补充。
}

\item
%% number: 10
%% paperName: 2013年江苏高考第 10 题 8 分
%% typeId: 4
%% type: 实验题
%% chapter: 电路
%% point: 电路实验
%% method: 无
%% score: 8
%% degree: 650
%% duplicateId: 0
%% body:
为探究小灯泡的电功率$P$和电压$U$的关系，小明测量小灯泡的电压$U$和电流$I$，利用$P=UI$得到电功率。实验所使用的小灯泡规格为“$3.0\UV$ $1.8\UW$”，电源为$12\UV$的电池，滑动变阻器的最大阻值为$10\UO$。
\begin{enumerate}
	\item
	准备使用的实物电路如题图1所示。请将滑动变阻器接入电路的正确位置。（用笔画线代替导线）
	\item
	现有$10\UO$、$20\UO$和$50\UO$的定值电阻，电路中的电阻$R_{1}$应选\tkanswer[1.2cm]{10}$\UO$的定值电阻。
	\item
	测量结束后，应先断开开关，拆除\tkanswer[1.2cm]{电池}两端的导线，再拆除其他导线，最后整理好器材。
	\item
	小明处理数据后将$P$、$U^{2}$点在坐标纸上，并作出了一条直线，如题图2示。请指出图象中不恰当的地方。
\end{enumerate}
\twopicture[numstyle=arabic, labelA=2013江苏10a, labelB=2013江苏10b, align=right, wA=4.8cm, wB=3.4cm]
{figs/10a.png}
{figs/10b.png}
%% answer:
\jdanswer{
\begin{enumerate}
	\item
	滑动变阻器采用限流式接法（将“一上一下”两个接线柱串联接入电路），连线见原卷答案图（原卷及材料未提供答案连线图，待补充）
	\item
	$10\UO$
	\item
	电池
	\item
	图线不应画成直线（小灯泡电阻随温度升高而增大，$P$ 与 $U^{2}$ 不成正比，应画成曲线），且横坐标 $U^{2}$ 的标度过大
\end{enumerate}
}
%% memo:
\memoanswer{
本题原卷及材料中未提供详解，待补充。
}

\item
%% number: 11
%% paperName: 2013年江苏高考第 11 题 19 分
%% typeId: 4
%% type: 实验题
%% chapter: 直线运动
%% point: 自由落体运动
%% method: 无
%% score: 19
%% degree: 450
%% duplicateId: 0
%% body:
某兴趣小组利用自由落体运动测定重力加速度，实验装置如图所示。倾斜的球槽中放有若干个小铁球，闭合开关K，电磁铁吸住第1个小球。手动敲击弹性金属片M，M与触头瞬间分开，第1个小球开始下落，M迅速恢复，电磁铁又吸住第2个小球。当第1个小球撞击M时，M与触头分开，第2个小球开始下落$\ldots\ldots$。这样，就可测出多个小球下落的总时间。
\begin{enumerate}
	\item
	在实验中，下列做法正确的有\tkanswer[1.2cm]{BD}。

	（A）电路中的电源只能选用交流电源

	（B）实验前应将M调整到电磁铁的正下方

	（C）用直尺测量电磁铁下端到M的竖直距离作为小球下落的高度

	（D）手动敲击M的同时按下秒表开始计时
	\item
	实验测得小球下落的高度$H=1.980\Um$，10个小球下落的总时间$T=6.5\Us$。可求出重力加速度$g=$\tkanswer[1.6cm]{9.4}$\Umsq$。（结果保留两位有效数字）
	\item
	在不增加实验器材的情况下，请提出减小实验误差的两个办法。
	\item
	某同学考虑到电磁铁在每次断电后需要时间$\Delta t$磁性才消失，因此，每个小球的实际下落时间与它的测量时间相差$\Delta t$，这导致实验误差。为此，他分别取高度$H_{1}$和$H_{2}$，测量$n$个小球下落的总时间$T_{1}$和$T_{2}$。他是否可以利用这两组数据消除$\Delta t$对实验结果的影响？请推导说明。
\end{enumerate}
\onepicture[w=4.5cm, align=right]{figs/11.png}
%% answer:
\jdanswer{
\begin{enumerate}
	\item
	BD
	\item
	9.4
	\item
	增大下落高度 $H$；增加每次下落的小球个数 $n$（多次测量取平均值）
	\item
	由 $H_{1}=\frac{1}{2}g(\frac{T_{1}}{n}-\Delta t)^{2}$ 和 $H_{2}=\frac{1}{2}g(\frac{T_{2}}{n}-\Delta t)^{2}$，可得 $g=\frac{2n^{2}(\sqrt{H_{1}}-\sqrt{H_{2}})^{2}}{(T_{1}-T_{2})^{2}}$，因此可以消去 $\Delta t$ 的影响。
\end{enumerate}
}
%% memo:
\memoanswer{
本题原卷及材料中未提供详解，待补充。
}

\item
%% number: 12
%% paperName: 2013年江苏高考第 12 题 4 分
%% typeId: 6
%% type: 计算题
%% chapter: 运动和力
%% point: 动量守恒定律
%% method: 无
%% score: 4
%% degree: 650
%% duplicateId: 0
%% body:
\textbf{A．}（选修模块 3-3）

如图所示，一定质量的理想气体从状态A依次经过状态B、C和D后再回到状态A。其中，A$\to$B和C$\to$D为等温过程，B$\to$C和D$\to$A为绝热过程（气体与外界无热量交换）。这就是著名的“卡诺循环”。

\onepicture[w=4.2cm, align=right]{figs/12a.png}

\begin{enumerate}
	\item
	该循环过程中，下列说法正确的是\xzanswer{C}

	\fourchoices[answer=C]
	{A$\to$B过程中，外界对气体做功}
	{B$\to$C过程中，气体分子的平均动能增大}
	{C$\to$D过程中，单位时间内碰撞单位面积器壁的分子数增多}
	{D$\to$A过程中，气体分子的速率分布曲线不发生变化}
	\item
	该循环过程中，内能减小的过程是\tkanswer[2.0cm]{B$\to$C}（选填“A$\to$B”、“B$\to$C”、“C$\to$D”或“D$\to$A”）。若气体在A$\to$B过程中吸收$63\UkJ$的热量，在C$\to$D过程中放出$38\UkJ$的热量，则气体完成一次循环对外做的功为\tkanswer[1.2cm]{25}$\UkJ$。
	\item
	若该循环过程中的气体为$1\Umol$，气体在A状态时的体积为$10\UL$，在B状态时压强为A状态时的$\frac{2}{3}$。求气体在B状态时单位体积内的分子数。（已知阿伏加德罗常数$N_{A}=6.02\times10^{23}\Umoln$，计算结果保留一位有效数字）
\end{enumerate}

\textbf{B．}（选修模块 3-4）

\begin{enumerate}
	\item
	如题图所示的装置，弹簧振子的固有频率是$4\UHz$。现匀速转动把手，给弹簧振子以周期性的驱动力，测得弹簧振子振动达到稳定时的频率为$1\UHz$，则把手转动的频率为\xzanswer{A}

	\onepicture[w=3.3cm, align=right]{figs/12b.png}

	\fourchoices[answer=A]
	{$1\UHz$}
	{$3\UHz$}
	{$4\UHz$}
	{$5\UHz$}
	\item
	如题图所示，两艘飞船A、B沿同一直线同向飞行，相对地面的速度均为$v$（$v$接近光速$c$）。地面上测得它们相距为$L$，则A测得两飞船间的距离\tkanswer[1.2cm]{大于}（选填“大于”、“等于”或“小于”）$L$。当B向A发出一光信号，A测得该信号的速度为\tkanswer[1.6cm]{$c$}。

	\onepicture[w=5.0cm, align=right]{\includesvg{figs/12c}}
	\item
	题图为单反照相机取景器的示意图，ABCDE为五棱镜的一个截面，AB$\perp$BC。光线垂直AB射入，分别在CD和EA上发生反射，且两次反射的入射角相等，最后光线垂直BC射出。若两次反射都为全反射，则该五棱镜折射率的最小值是多少？（计算结果可用三角函数表示）

	\onepicture[w=4.5cm, align=right]{figs/12d.png}
\end{enumerate}

\textbf{C．}（选修模块 3-5）

\begin{enumerate}
	\item
	如果一个电子的德布罗意波长和一个中子的相等，则它们的\xzanswer{C}也相等。

	\fourchoices[answer=C]
	{速度}
	{动能}
	{动量}
	{总能量}
	\item
	根据玻尔原子结构理论，氦离子（\ce{He+}）的能级图如题图所示。电子处在$n=3$轨道上比处在$n=5$轨道上离氦核的距离\tkanswer[1.0cm]{近}（选填“近”或“远”）。当大量\ce{He+}处在$n=4$的激发态时，由于跃迁所发射的谱线有\tkanswer[1.0cm]{6}条。

	\onepicture[w=4.2cm, align=right]{figs/12e.png}
	\item
	如题图所示，进行太空行走的宇航员A和B的质量分别为$80\Ukg$和$100\Ukg$，他们携手远离空间站，相对空间站的速度为$0.1\Ums$。A将B向空间站方向轻推后，A的速度变为$0.2\Ums$，求此时B的速度大小和方向。

	\onepicture[w=6.0cm, align=right]{figs/12f.png}
\end{enumerate}
%% answer:
\jdanswer{
\begin{enumerate}
	\item
	A：（1）C；（2）B$\to$C，$25\UkJ$；（3）等温过程 $p_{A}V_{A}=p_{B}V_{B}$，单位体积内的分子数 $n=\frac{N_{A}}{V_{B}}$，解得 $n=\frac{N_{A}p_{B}}{p_{A}V_{A}}$，代入数据得 $n=4\times10^{25}\Umnc$。
	\item
	B：（1）A；（2）大于，$c$（或光速）；（3）由题意知，入射角 $\alpha=22.5^{\circ}$，则折射率最小值 $n=\frac{1}{\sin22.5^{\circ}}$。
	\item
	C：（1）C；（2）近，6；（3）根据动量守恒 $(m_{A}+m_{B})v_{0}=m_{A}v_{A}+m_{B}v_{B}$，解得 $v_{B}=0.02\Ums$，方向为离开空间站。
\end{enumerate}
}
%% memo:
\memoanswer{
本题原卷及材料中未提供详解，待补充。
}

\item
%% number: 13
%% paperName: 2013年江苏高考第 13 题 15 分
%% typeId: 6
%% type: 计算题
%% chapter: 电磁感应
%% point: 电磁感应电路分析
%% method: 无
%% score: 15
%% degree: 650
%% duplicateId: 0
%% body:
如图所示，匀强磁场中有一矩形闭合线圈abcd，线圈平面与磁场垂直。已知线圈的匝数$N=100$，边长$ab=1.0\Um$、$bc=0.5\Um$，电阻$r=2\UO$。磁感应强度$B$在$0\sim1\Us$内从零均匀变化到$0.2\UT$。在$1\sim5\Us$内从$0.2\UT$均匀变化到$-0.2\UT$，取垂直纸面向里为磁场的正方向。求：
\begin{enumerate}
	\item
	$0.5\Us$时线圈内感应电动势的大小$E$和感应电流的方向；
	\item
	在$1\sim5\Us$内通过线圈的电荷量$q$；
	\item
	在$0\sim5\Us$内线圈产生的焦耳热$Q$。
\end{enumerate}
\onepicture[w=5.0cm, align=right]{figs/13.png}
%% answer:
\jdanswer{
\begin{enumerate}
	\item
	$10\UV$，感应电流的方向 $a\to d\to c\to b\to a$
	\item
	$10\UC$
	\item
	$100\UJ$
\end{enumerate}
}
%% memo:
\memoanswer{
（1）感应电动势 $E_{1}=N\frac{\Delta\Phi_{1}}{\Delta t_{1}}$，磁通量的变化 $\Delta\Phi_{1}=\Delta B_{1}S$。

解得 $E_{1}=N\frac{\Delta B_{1}S}{\Delta t_{1}}$，代入数据得 $E_{1}=10\UV$，感应电流的方向 $a\to d\to c\to b\to a$。

（2）同理可得 $E_{2}=N\frac{\Delta B_{2}S}{\Delta t_{2}}$，感应电流 $I_{2}=\frac{E_{2}}{r}$，电量 $q=I_{2}\Delta t_{2}$。

解得 $q=N\frac{\Delta B_{2}S}{r}$，代入数据得 $q=10\UC$。

（3）$0\sim1\Us$内的焦耳热 $Q_{1}=I_{1}^{2}r\Delta t_{1}$，且 $I_{1}=\frac{E_{1}}{r}$；$1\sim5\Us$内的焦耳热 $Q_{2}=I_{2}^{2}r\Delta t_{2}$。

由 $Q=Q_{1}+Q_{2}$，代入数据得 $Q=100\UJ$。
}

\item
%% number: 14
%% paperName: 2013年江苏高考第 14 题 16 分
%% typeId: 6
%% type: 计算题
%% chapter: 运动和力
%% point: 牛顿定律的应用
%% method: 无
%% score: 16
%% degree: 650
%% duplicateId: 0
%% body:
如图所示，将小砝码置于桌面上的薄纸板上，用水平向右的拉力将纸板迅速抽出，砝码的移动很小，几乎观察不到，这就是大家熟悉的惯性演示实验。若砝码和纸板的质量分别为$m_{1}$和$m_{2}$，各接触面间的动摩擦因数均为$\mu$。重力加速度为$g$。
\begin{enumerate}
	\item
	当纸板相对砝码运动时，求纸板所受摩擦力的大小；
	\item
	要使纸板相对砝码运动，求所需拉力的大小；
	\item
	本实验中，$m_{1}=0.5\Ukg$，$m_{2}=0.1\Ukg$，$\mu=0.2$，砝码与纸板左端的距离$d=0.1\Um$，取$g=10\Umsq$。若砝码移动的距离超过$l=0.002\Um$，人眼就能感知。为确保实验成功，纸板所需的拉力至少多大？
\end{enumerate}
\onepicture[w=6.0cm, align=right]{figs/14.png}
%% answer:
\jdanswer{
\begin{enumerate}
	\item
	$\mu(2m_{1}+m_{2})g$
	\item
	$F>2\mu(m_{1}+m_{2})g$
	\item
	$22.4\UN$
\end{enumerate}
}
%% memo:
\memoanswer{
（1）砝码对纸板的摩擦力 $f_{1}=\mu m_{1}g$，桌面对纸板的摩擦力 $f_{2}=\mu(m_{1}+m_{2})g$，$f=f_{1}+f_{2}$。解得 $f=\mu(2m_{1}+m_{2})g$。

（2）设砝码的加速度为 $a_{1}$，纸板的加速度为 $a_{2}$，则 $f_{1}=m_{1}a_{1}$，$F-f_{1}-f_{2}=m_{2}a_{2}$，发生相对运动 $a_{2}>a_{1}$。解得 $F>2\mu(m_{1}+m_{2})g$。

（3）纸板抽出前，砝码运动的距离 $x_{1}=\frac{1}{2}a_{1}t_{1}^{2}$，纸板运动的距离 $d+x_{1}=\frac{1}{2}a_{2}t_{1}^{2}$。纸板抽出后，砝码在桌面上运动的距离 $x_{2}=\frac{1}{2}a_{3}t_{2}^{2}$，$l=x_{1}+x_{2}$。由题意知 $a_{1}=a_{3}$，$a_{1}t_{1}=a_{3}t_{2}$，解得 $F=2\mu[m_{1}+(1+\frac{d}{l})m_{2}]g$。代入数据得 $F=22.4\UN$。
}

\item
%% number: 15
%% paperName: 2013年江苏高考第 15 题 16 分
%% typeId: 6
%% type: 计算题
%% chapter: 磁场
%% point: 带电粒子在磁场中的运动
%% method: 无
%% score: 16
%% degree: 350
%% duplicateId: 0
%% body:
在科学研究中，可以通过施加适当的电场和磁场来实现对带电粒子运动的控制。如题图1所示的$xOy$平面处于匀强电场和匀强磁场中，电场强度$E$和磁感应强度$B$随时间$t$作周期性变化的图象如题图2所示。$x$轴正方向为$E$的正方向，垂直纸面向里为$B$的正方向。在坐标原点$O$有一粒子P，其质量和电荷量分别为$m$和$+q$。不计重力。在$t=\frac{\tau}{2}$时刻释放P，它恰能沿一定轨道做往复运动。
\begin{enumerate}
	\item
	求P在磁场中运动时速度的大小$v_{0}$；
	\item
	求$B_{0}$应满足的关系；
	\item
	在$t_{0}$（$0<t_{0}<\frac{\tau}{2}$）时刻释放P，求P速度为零时的坐标。
\end{enumerate}
\twopicture[numstyle=arabic, labelA=2013江苏15a, labelB=2013江苏15b, align=right, wA=3.2cm, wB=3.0cm]
{figs/15a.png}
{figs/15b.png}
%% answer:
\jdanswer{
\begin{enumerate}
	\item
	$v_{0}=\frac{qE_{0}\tau}{2m}$
	\item
	$B_{0}=\frac{(2n-1)\pi m}{q\tau}$（$n=1$，2，3$\cdots$）
	\item
	纵坐标为 $\frac{2E_{0}[k(\tau-2t_{0})+t_{0}]}{B_{0}}$ 或 $\frac{2E_{0}k(\tau-2t_{0})}{B_{0}}$（$k=1$，2，3$\cdots$）
\end{enumerate}
}
%% memo:
\memoanswer{
（1）$\frac{\tau}{2}\sim\tau$ 作匀加速直线运动，$\tau\sim2\tau$ 作匀速圆周运动。电场力 $F=qE_{0}$，加速度 $a=\frac{F}{m}$，速度 $v_{0}=at$，且 $t=\frac{\tau}{2}$。解得 $v_{0}=\frac{qE_{0}\tau}{2m}$。

（2）只有当 $t=2\tau$ 时，P在磁场中作圆周运动结束并开始沿$x$轴负方向运动，才能沿一定轨道作往复运动，如图所示。设P在磁场中做圆周运动的周期为$T$。则 $(n-\frac{1}{2})T=\tau$（$n=1,2,3\cdots$）。匀速圆周运动 $qvB_{0}=m\frac{v^{2}}{r}$，$T=\frac{2\pi r}{v}$。解得 $B_{0}=\frac{(2n-1)\pi m}{q\tau}$（$n=1,2,3\cdots$）。

\onepicture[w=3.2cm]{figs/15_an1.png}

（3）在$t_{0}$时刻释放，P在电场中加速时间为$\tau-t_{0}$。在磁场中匀速圆周运动 $v_{1}=\frac{qE_{0}(\tau-t_{0})}{m}$，圆周运动的半径 $r_{1}=\frac{mv_{1}}{qB_{0}}$。解得 $r_{1}=\frac{E_{0}(\tau-t_{0})}{B_{0}}$。又经$(\tau-t_{0})$时间P减速为零后向右加速时间为$t_{0}$。P再进入磁场 $v_{2}=\frac{qE_{0}t_{0}}{m}$，圆周运动的半径 $r_{2}=\frac{mv_{2}}{qB_{0}}$。解得 $r_{2}=\frac{E_{0}t_{0}}{B_{0}}$。综上分析，速度为零时横坐标$x=0$，$y=\begin{cases}2[kr_{1}-(k-1)r_{2}]\\2k(r_{1}-r_{2})\end{cases}$（$k=1,2,3\cdots$）。

\onepicture[w=3.2cm]{figs/15_an2.png}

解得相应的纵坐标为 $\frac{2E_{0}[k(\tau-2t_{0})+t_{0}]}{B_{0}}$，$\frac{2E_{0}k(\tau-2t_{0})}{B_{0}}$（$k=1,2,3\cdots$）。
}

\end{enumerate}

\end{document}
